\documentclass[11pt]{article}
\usepackage[T1]{fontenc}
\usepackage{lmodern,amsmath,amssymb,amsthm,mathtools}
\usepackage[margin=1in]{geometry}
\usepackage{booktabs,longtable,array,enumitem,microtype,xcolor}
\usepackage{listings,fancyhdr}
\usepackage[hidelinks]{hyperref}
\usepackage{xurl}
\definecolor{navy}{RGB}{28,54,79}
\definecolor{pale}{RGB}{244,246,248}
\hypersetup{pdftitle={K215: automatic source certificates, concrete credit proofs, and the 214-state search},pdfauthor={Research contribution for Vladimir Reshetnikov's ProveIt}}
\pagestyle{fancy}\fancyhf{}\lhead{\small K215: proof simplification}\rhead{\small 21 October 2026}\cfoot{\thepage}
\setlength{\headheight}{14pt}
\setlength{\parskip}{5pt}\setlength{\parindent}{0pt}
\clubpenalty=10000\widowpenalty=10000\displaywidowpenalty=10000
\newcolumntype{L}[1]{>{\raggedright\arraybackslash}p{#1}}
\setlist{itemsep=3pt,topsep=5pt}
\lstset{basicstyle=\ttfamily\footnotesize,backgroundcolor=\color{pale},frame=single,rulecolor=\color{pale},breaklines=true,columns=fullflexible,keepspaces=true,showstringspaces=false}
\newtheorem{theorem}{Theorem}[section]
\newtheorem{lemma}[theorem]{Lemma}
\newtheorem{proposition}[theorem]{Proposition}
\theoremstyle{definition}\newtheorem{definition}[theorem]{Definition}
\newcommand{\Con}{\operatorname{Con}}
\newcommand{\Halt}{\operatorname{Halt}}
\newcommand{\pair}[2]{\mathsf P(#1,#2)}
\newcommand{\Drop}{\mathsf{DROP}}
\newcommand{\Nip}{\mathsf{NIP}}
\newcommand{\MP}{\mathsf{MP}}
\newcommand{\AX}{\mathsf{AX}}
\newcommand{\AutoStatus}{The final source theorem passed kernel checking. All thirty selected project modules were compiled from source in the isolated build; successful modules were hash-bound across the documented resource retries. The final printed axioms are exactly \texttt{propext}, \texttt{Classical.choice}, and \texttt{Quot.sound}.}
\newcommand{\RocqStatus}{All twelve selected project modules compiled freshly, and ordinary recursive \texttt{coqchk} passed without \texttt{-admit} or \texttt{-norec}.}
% Filled from verified receipts before final delivery.
\newcommand{\RocqMeasurements}{All twelve compiler children took 65.99 seconds in total, with peak resident memory 816,160 KiB. The concrete source module took 15.35 seconds; its forward \texttt{vm\_compute} transaction reported 7.006 seconds. Recursive \texttt{coqchk} took 543.73 seconds and peaked at 338,332 KiB. All five new printed assumption reports were closed. These are local measurements of the selected source closure, not a full-project speedup comparison.}
\newcommand{\AutoMeasurements}{The thirty accepted compiler invocations took 578.41 seconds in total; the successful process records, including emitter replays, sum to 588.20 seconds. Peak accepted resident memory was 3,954,196 KiB. Two staged-check resource failures consumed a separately recorded 162.09 seconds. These are sums of recorded process times across the isolated development build, not the wall time of a single uninterrupted clean command. The final cache regenerates byte for byte; the supplied default fresh-build driver uses the accepted final sources and a 4,000 MiB cap.}
\newcommand{\SearchMeasurements}{The broad run tested 7,324,620 local replacement actions across 21 released tables and performed 8,475,730 guarded case tests. It found no complete state elimination. A targeted provider search tested 731,000 provider actions, retained 1,812 by its preservation criterion, and then tested 1,558,320 root actions without a successful repair.}


\begin{document}
\begin{center}
{\Large\bfseries\color{navy} K215: automatic source certificates,\\[3pt]
concrete credit proofs, and the 214-state search}\par
\vspace{8pt}
{\large A research and integration report for ProveIt}\par
21 October 2026
\end{center}

\begin{abstract}
The current 215-state machine already has checked Lean and Rocq correctness
endpoints. The useful next step is therefore to preserve that calibration and
remove repeated source-correspondence and enumeration arguments. This report
reviews all seven pending proposal archives, supplies a new Rocq-compatible
normalized source checker with an exact all-register theorem, implements
automatic correspondence inference in Lean from the two literal decoder trees,
and proves the concrete zero-credit scheduler theorem for the existing ZC215
codec. The first construction reduces the source relation from 2,240 entries to
1,426. The second removes a separately maintained correspondence input by
synthesizing a cache and checking it against the source trees. The third discharges
the actual codec, parser, monotonicity, and injectivity assumptions behind the
previous abstract credit argument. A separate predecessor-specific state
elimination search and a targeted scanner repair search did not construct a
verified 214-state machine. Their finite scope is recorded explicitly.
\end{abstract}

\section{Results, scope, and recommendation}

The reviewed repository is pinned to commit
\texttt{53269bcab1efc0f82d3022645bd7cf0c22d84959}.
Unless stated otherwise, paths below are relative to
\texttt{Computability/BusyBeaver/Research/ZFC/}.
All seven ZIP archives then present in \texttt{docs/incoming/} were inspected;
their filenames, hashes, and individual contributions are recorded in the
accompanying incoming-proposal inventory. The newest archive is the one whose
filename ends in \texttt{(2).zip}. The main candidate and the established
endpoints are inherited work, not new results of this report.\cite{verification,incoming}

\begin{center}
\begin{tabular}{L{.24\textwidth}L{.39\textwidth}L{.27\textwidth}}
\toprule
Component & Contribution here & Exact boundary\\
\midrule
Normalized Rocq source proof & 1,426-entry checker and all-register equivalence
for literal \texttt{p215} and \texttt{dagPerm} & Coq 8.18.0 build; full binary/ZFC
closure and Rocq 9.1 integration are separate\\[3pt]
Automatic Lean certificate & Infer and cache paths and successors from the two
decoder trees; validate with the proved normalized checker & Source equivalence;
no machine-state reduction\\[3pt]
Concrete ZC215 credit & Exact Cantor codec, parser inverses, lowering injection,
and all-prefix credit in one closed development & Scheduler theorem; not a new
ZC215 binary endpoint\\[3pt]
State reduction & Predecessor-dependent elimination across 21 released tables,
plus scanner repair & No accepted 214-state construction; no lower bound\\
\bottomrule
\end{tabular}
\end{center}

The practical recommendation has three parts. First, retain the already
certified K215 as the primary endpoint and adopt proof refactors behind the
same public statements. Second, use the existing conditional-marker CM215
when the aim is the shortest additional formal correctness proof: the incoming
packets already contain its exact relative halting bridge. Third, use ZC215
when simplifying the enumerator's mathematical explanation is the priority;
its tag-lowering injection gives an especially direct protected-stack argument.
These are different measures of simplicity. A short proof relative to K215
still inherits K215's foundational and compiler development.

\textbf{214-state answer.} No verified 214-state or smaller machine was
constructed in this investigation. No rejected search output is offered as
a candidate satisfying the consistency specification. The useful positive
deliverables are the proof components and the reproducible search framework.

\subsection{Meaning of the target}

The desired statement is
\[
  \mathsf{ZFC}\vdash
  \bigl(\neg\Halt(M,\mathrm{blank})\ \longleftrightarrow\ \Con(\mathsf{ZFC})\bigr),
\]
where consistency is syntactic nonexistence of a finite refutation in the
specified proof system. This is an equivalence theorem, not an unconditional
claim that either side holds. Its elementary machine and source transformations
can be formalized in ZFC. A checked external Lean or Rocq theorem about the
corresponding predicates must also be distinguished from a proof object
encoded inside a formalization of ZFC or PA itself. The repository documents
an internal PA endpoint for the earlier 240-state machine; the current K215
external endpoint does not automatically settle that separate milestone.
\cite{verification}

All machine counts use the conventional binary, one-tape, two-way infinite
tape model with L/R moves, blank input, and initialization included. The
separate halt is excluded from the number of working states. A 215-state
table therefore has 430 working transition rows. Merely using a large ambient
type for labels does not add reachable working states; conversely, renaming
states does not remove one.

\section{What the existing proposals already establish}

The current repository contains substantial simplification work. Its phase-four
record reports a 185-module Lean closure with ten freshly rebuilt affected
modules and 175 precisely bound reused objects. It shares implication lemmas,
removes a redundant segment-halting induction, and proves enumerator blocks
through symbolic execution. The three highlighted files decrease from 1,814
to 1,606 lines. Those are existing improvements and are not credited to this
packet. The underlying full endpoint has a separate fresh-build history.
\cite{phase4}

The incoming proposals also overlap considerably. A review that reads only the
main candidate's README misses several already-compiled follow-up results.
The following inventory records the main distinctions; each entry is an
archive-level claim, not an assertion that its entire inherited endpoint was
freshly rebuilt during this work.

\begingroup\small
\begin{longtable}{L{.34\textwidth}L{.57\textwidth}}
\toprule
Archive stem, all dated 2026-10-21 & Already supplied contribution\\
\midrule\endhead
\texttt{K215\_Proof\_Simplification (2)} & Exact CM215 dense/Core bridge;
ordinary zero-offset source equivalence; Coq certificate identity and mirror
checker; a final full-ZFC adapter with an explicit integration boundary.\\
\texttt{K215\_Proof\_Simplification} & Compiled 1,426-entry normalized Lean
source equivalence; compiled CM215 transfer; generic positive stuttering.\\
\texttt{K215\_Proof\_Transport} & Exact CM215 Core bridge, positive-block
theorems, source-graph measurements, and bounded reduction searches.\\
\texttt{Structural\_Proof\_\allowbreak Simplification} & Equal canonical weighted test
graphs established by finite analysis; structural and semantic proposals.
Its canonical certificate is not accepted by the unchanged raw-path checker.\\
\texttt{ZFC215\_Proof\_Simplification} & Plain one-sided Coq source proof,
compact paths, and a marker-variant development with stated adapter limits.\\
\texttt{ZFC\_K215\_Proof\_Transport} & Standalone all-tape CM215 proof,
shared-tree and compact-path pilots, with the remaining representation
adapter identified.\\
\texttt{zfc-fm215-proof-compression} & Four-row FM215 variant, standalone
all-tape equivalence, decoder-context compression, and a restricted
tail-quantifier countermodel.\\
\bottomrule
\end{longtable}
\endgroup

Two consequences guide the present work. Recommending mirror symmetry,
zero-offset specialization, or shared trees alone would repeat existing
proposals. More substantially, the 1,426-entry Lean certificate is already
checked, so its case-count reduction is not a new discovery here. The new
work is to bring the Rocq lane to that reduced relation and to remove an
explicit correspondence file from the Lean-facing interface. The concrete
scheduler proof similarly formalizes an existing mathematical idea while
closing hypotheses that the old abstract proof left to its caller.

\section{Normalized source correspondence in the Rocq lane}

\subsection{Why 814 entries can be removed}

The path machine reads a finite Boolean pointer as if it were followed by
infinitely many zero bits. Define
\[
 P\simeq Q\quad\Longleftrightarrow\quad
 \forall i\in\mathbb N,\ \operatorname{bit}(P,i)=\operatorname{bit}(Q,i),
\]
where absent bits are zero. For example, \((1,0)\simeq(1,0,0,0)\).
After it finds a primitive leaf, the decoder executes using the actual
consumed path. Padding is therefore a representational artifact, not an
observable difference in the register computation.

The old certificate has 2,240 path pairs. Normalizing the trailing-zero aliases
leaves 1,426 relevant pairs, a reduction of
\[
  2240-1426=814,\qquad \frac{814}{2240}\approx 36.34\%.
\]
There are 1,425 decrement checkpoints and one halt checkpoint. Both outcomes
of every decrement still require validation, as does the initial segment.
The reduction does not throw away a branch because it was absent from a
bounded execution trace.

\begin{lemma}[Path aliases preserve halting]
For any source program $p$, register file $R$, and pointers $P\simeq Q$,
the executions from $(P,R)$ and $(Q,R)$ have the same halting predicate.
The finite halting-witness transport is already available in the inherited
path semantics.
\end{lemma}

\begin{proof}
The decoder observes identical bits at every position, so it selects the same
primitive instruction and consumed depth. The zero-extended prefixes of that
depth agree. Consequently the next register update, next pointer, and literal
halt observation agree. Transport a finite run by induction. The formal
development reuses \texttt{RMEquiv.hn\_ipeq}; it does not introduce a second
path semantics or reinterpret a failed decoder as successful halting.
\end{proof}

The executable Boolean comparison \texttt{ibeq} checks this relation without
constructing an infinite bit sequence. It compares common prefixes and accepts
an unmatched suffix exactly when all its bits are false. Its soundness theorem
is closed under the global context.

\subsection{Segments and the reduced checker}

Between decrement tests, the source executes only increments and structural
break/continue instructions. Such a segment does not inspect a register.
It has a finite increment-count vector and a final control pointer. The
inherited evaluator \texttt{seg} computes that information with a fuel bound;
fuel exhaustion is rejected.

For a certificate entry $(P_1,P_2,j,k)$, the new checker verifies either that
both leaves are halt, or that both decrement the same register and:
\begin{enumerate}
\item the nonzero branches have equal increment counts and end at pointers
equivalent under $\simeq$ to the pair indexed by $j$;
\item the zero branches satisfy the corresponding condition for $k$.
\end{enumerate}
The existing source programs, segment evaluator, register updates, and halting
predicate remain the inputs to this test. Only the endpoint comparison in the
certificate checker changes from literal finite-list equality to $\simeq$.

\begin{theorem}[One checked normalized certificate suffices]
If the initial segment condition and every normalized entry condition succeed,
then for every register file $R$,
\[
 \operatorname{rm\_halts}(p_1,([],R))
 \quad\Longleftrightarrow\quad
 \operatorname{rm\_halts}(p_2,([],R)).
\]
\end{theorem}

\begin{proof}
For the forward implication, induct on the length of a finite halting run.
A related halt pair ends the argument. At a decrement pair, equal register
files choose the same branch. That decrement consumes a source step, so the
induction parameter strictly decreases. The two following segments produce
the same register increments. Transport their final pointers to the indexed
certificate pointers by the alias lemma, and apply the induction hypothesis.
The initial segment is handled in the same way, without needing it to have
positive length.

For the reverse implication, swap the two pointer components of every entry.
The increment-count equality is symmetric, and the two endpoint alias checks
exchange positions. A generic mirror theorem establishes acceptance of the
mirrored certificate from the already-proved forward acceptance. Apply the
same soundness theorem. No second large computed certificate check is needed.
\end{proof}

\subsection{Concrete checked theorem and costs}

The delivered \texttt{Source215Normalized.v} checks the literal inherited
\texttt{Prog215.v} and \texttt{DagPerm.v} and proves:
\begin{lstlisting}
Theorem p215_dagPerm_normalized R :
  rm_halts p215 ([], R) <-> rm_halts dagPerm ([], R).
\end{lstlisting}
The Boolean certificate theorem and this endpoint print closed assumption
reports. \RocqStatus

The active data source is 109,115 bytes. The incoming ordinary 2,240-entry
counterpart is 171,678 bytes, so the like-for-like data saving is 62,563 bytes,
approximately 36.44\%. This comparison should not be confused with the much
larger saving obtained earlier by replacing expanded Boolean path lists with
binary-coded paths. The new generic checker is an additional reusable proof
component and must be counted when measuring the total development.

The new reusable checker is 196 lines (8,032 bytes), and the concrete source
instance is 24 lines (901 bytes). Relative to production's two separately
computed 2,240-entry directions, concrete entry checks decrease from 4,480 to
1,426. The mirror technique already appeared in the earlier one-sided proposal;
only the normalization-aware Rocq proof and instance are new here.

\RocqMeasurements

The reproducible build uses Coq 8.18.0. A small, proved compatibility layer
provides renamed standard-library facts required by the imported Rocq-era
source; the original source and the exact build transformation are retained.
This is not a claim that the complete 237-source production endpoint has
been rebuilt under Rocq 9.1.0. The supplied integration wrapper preserves the
old source theorem's public role, but production adoption should be followed
by that repository's normal full endpoint and assumption audit.

\section{Automatic Lean correspondence inference}

\subsection{Removing a separately maintained input}

The normalized Lean proposal still ships an explicit 1,426-entry table of
paired pointers and successor indices. The structural proposal observes that
the two trees' test leaves correspond by their order. We turn that observation
into an executable certificate generator and retain the existing proved
checker as the acceptance boundary.

The generator takes only the two decoder trees. It performs the following
finite operations:
\begin{enumerate}
\item Traverse each tree from left to right and list the full consumed paths
of its decrement and halt leaves.
\item Pair these lists in order. Encode each path by a leading sentinel one
followed by its bits.
\item For each source decrement, compute its zero and nonzero continuation
segments in the first tree. Determine the resulting checkpoint's rank in
that tree. Use those ranks as the proposed successor indices.
\item Build a balanced certificate and run the normalized segment checker
against \emph{both} actual trees.
\end{enumerate}

The rank computation follows the pointer down the decoder tree, adding the
number of checkpoints in a left subtree whenever it takes a right edge. It
therefore avoids a linear search through every checkpoint path. As in the
source semantics, exhausted pointer lists are interpreted with trailing zeros.

The generator is deliberately a heuristic rather than a theorem that every
equivalent pair of source programs has this particular ordering. A pair of
equivalent programs with a different checkpoint order can be rejected. That
incompleteness is harmless: acceptance is established by the already-proved
checker, not by an assumed correspondence property of the generator.

\begin{theorem}[Sound automatic inference]
Let $C=\operatorname{synthesize}(T_1,T_2)$. If the normalized initial and entry
checks succeed for $C$, and $T_i$ are the actual decoder trees of $p_i$, then
the two source programs have equivalent halting predicates from the empty
pointer for every register file.
\end{theorem}

\begin{proof}
Instantiate the normalized checker's soundness theorem with $C$. Nothing
about how $C$ was obtained is required. The proof is valid even if a generator
branch uses a fallback index, construction fuel is insufficient, or the two
path lists have different lengths: those situations must still survive every
semantic acceptance condition. In particular, an empty generated relation
does not pass the required initial relation check. The theorem consequently
adds no trusted Python computation, hash-based equality, or generator-correctness
axiom.
\end{proof}

This formulation is smaller than a full verified compiler from the source
language to a new graph semantics. It uses the structural proposal's insight
without requiring a second operational semantics and adequacy proof. It also
has a precise limit: it removes the \emph{manually maintained} correspondence
input, not the mathematical relation or the checking work. The executable
generator computes the relation, emits literal data, and submits that data to
the established kernel checker.

\subsection{Actual K215 instantiation}

The generic inference module is 88 lines. A small executable Lean emitter
computes the relation from the existing exact shared trees and partitions it
into sixteen literal checking blocks. The emitted data is a reproducible cache,
not a trusted proof. The final theorem uses the same public all-register target:
\begin{lstlisting}
theorem source_halts_iff (R : Registers) :
  Halts p215 <[], R> <-> Halts pathDagPerm <[], R>
\end{lstlisting}
Here the angle brackets in the display stand for Lean's configuration
constructor. The delivered file contains the actual Lean syntax.

The acceptance computations are expressed through \texttt{decide +kernel}.
The development does not use \texttt{native\_decide}, an admitted theorem,
or a finite-run assertion in place of the all-register argument. An initial
executable evaluation independently returned 1,426 rows, a successful initial
check, and successful checks for all entries. Executable evaluation alone is
not presented as a proof of the concrete result. \AutoStatus

\AutoMeasurements

The current repository's \texttt{RMBisim.lean}, including its phase-four
successful-prefix cancellation lemma, is used in the fresh build. The eight
inherited source-semantics/DAG dependencies are bound to the current commit by
Git blob identity. The source programs' literal node table and decoder-tree
bindings are also audited. Shared trees and the normalized checker are credited
to the incoming proposal from which they are reused.

\subsection{Engineering tradeoff}

The first implementation made each checking block reduce the inferred relation
definition directly. Its executable evaluation accepted all 1,426 entries, and
the first three blocks passed kernel checking, but block four exceeded the
configured 3,000 MiB memory limit. That failure is retained in the experimental
receipt; it is not a mathematical rejection of the certificate.

The delivered production-style experiment therefore stages the computation.
Lean executes the maintained inference algorithm once and writes a literal
balanced certificate. The kernel then validates every entry against the two
actual source trees. There is deliberately no theorem equating that cache with
the generator's output: such an equality is unnecessary for semantic soundness
and would repeat the expensive computation. The exact source-equivalence
theorem depends on checker acceptance, irrespective of who generated the data.
Later checking blocks also use the incoming proposal's quarter-block batching
to bound kernel memory. Resource retries are recorded separately from accepted
compilations. The build receipt records reproducibility of the delivered cache, which
is useful evidence about engineering identity but is not a logical premise.

A smaller source file is not automatically a faster proof. Synthesis performs
additional tree traversals, and evaluating a definition again in multiple
checking blocks can be expensive. The recorded times are measurements of this
build, not a benchmark against a fresh complete baseline under identical load.
The report makes no speedup claim without such a comparison.

\subsection{One inferred correspondence for both assistants}

The same newly inferred 1,426 tuples also pass the Rocq normalized checker.
The additional \texttt{rocq\_auto/} component renders the emitter's literal
\texttt{entries.json} through the existing sentinel-bit path codec and proves
the exact all-register \texttt{p215}/\texttt{dagPerm} theorem again. This is a
useful cross-assistant interface: maintain one inference algorithm, then let
the independently established semantic checkers validate its output against
each assistant's literal programs. Neither checker assumes the correctness
of the generator or of the other assistant.

The two new Coq modules compiled in 11.798 and 19.429 seconds; the acceptance
transaction reported 7.378 seconds. All three printed assumption reports are
closed. This additional check imports ten source-and-object-hash-bound
dependencies from the already accepted normalized build. It is an incremental
two-module validation, not another fresh dependency build or recursive
\texttt{coqchk} run. The first normalized source theorem retains the separate
successful recursive audit reported above.

The automatically inferred Coq data is 110,668 bytes, slightly larger than
the earlier 109,115-byte normalized data because it retains full consumed
paths. The benefit is a single maintained inference procedure and two semantic
validations, not a claim that every generated representation is smaller.

\section{A concrete, closed prefix-credit theorem for ZC215}

\subsection{Literal codec}

ZC215 is an existing 215-state alternative. Its important source change is that
its scheduler processes raw candidates $0,1,2,\ldots$, with DROP assigned tag
zero and NIP tag one. The old mathematical tag-lowering proof is already in
the repository. The older \texttt{PrefixCredit.v} draft abstracted the lowering
map and its properties. Our new development proves those properties for the
actual codec and a total executable parser.\cite{zc215}

Write Cantor pairing in the source's coordinate convention as
\[
 \pair ab=\frac{(a+b)(a+b+1)}2+b.
\]
The command fields are:
\[
\begin{array}{c|c}
\text{command}&\text{fields}\\\hline
\Drop &[0]\\
\Nip &[1]\\
\MP(x)&[2,x]\\
\AX(c;x,y)&[c+3,y,x]
\end{array}
\]
Let $\operatorname{emit}(k,r)$ nest those fields through pairing, with final
residual $r$. Define
\[
 \operatorname{code}([])=0,\qquad
 \operatorname{code}(k::w)=1+\operatorname{emit}(k,\operatorname{code}(w)).
\]
The raw candidate of a nonempty word is its code minus one. This final shift
must be retained: raw candidate zero represents the actual command word
$[\Drop]$, and raw candidate two represents $[\Drop,\Drop]$.

\begin{theorem}[Concrete parser inverses]
The delivered executable parser $\operatorname{word}$ satisfies
\[
 \operatorname{code}(\operatorname{word}(n))=n+1
\]
for every $n$, and
\[
 \operatorname{word}(\operatorname{code}(w)-1)=w
\]
for every nonempty word $w$.
\end{theorem}

The formal proof constructs Cantor unpairing, proves the required inverse and
coordinate inequalities, and recurses on a strictly smaller residual. Thus
the parser and the code injection are proved, not supplied as abstract
premises. The explicit residual convention prevents a common error: a final
zero command tag must not be mistaken for an absent command.

\subsection{Tag lowering is a well-founded injection}

An occurrence is a pair $(n,j)$ consisting of a raw candidate and a command
position. For a NIP occurrence, replace that NIP with DROP in
$\operatorname{word}(n)$, re-encode the resulting word as $m$, and set
\[
 L(n,j)=(m,j).
\]

\begin{lemma}[Three properties of lowering]
The image is a DROP occurrence, $m<n$, and $L$ is injective on NIP occurrences.
\end{lemma}

\begin{proof}
The first assertion follows from the parser inverse: decoding the new word's
code recovers that word, with DROP at the changed position. For strict
decrease, pairing is strictly increasing in either coordinate when the other
coordinate is fixed. Replacing tag one by tag zero therefore lowers the code
at the chosen command. Every preceding command encoder is strictly increasing
in its residual argument, so the decrease propagates to the entire raw code.

For injectivity, suppose two lowered occurrences are equal. Their candidate
and position agree. Decode that candidate and restore NIP at the specified
position. This recovers both original words. The code/parser inverse then
recovers the same original raw candidate. Hence both components of the
original occurrences agree.
\end{proof}

Notice that the image is in a \emph{strictly earlier candidate}, not merely
at an earlier position in the same candidate. This is what makes the argument
work uniformly at every prefix, including prefixes that stop immediately
before a NIP.

\subsection{Every canonical enumeration prefix has credit}

Let $\operatorname{prefix}(n,j)$ contain all occurrences in candidates
$0,\ldots,n-1$, followed by positions $0,\ldots,j-1$ of candidate $n$.
Let $D(n,j)$ and $N(n,j)$ be the numbers of DROP and NIP occurrences there.

\begin{theorem}[Nonnegative credit and strict pre-pop credit]
For all $n,j$,
\[
 N(n,j)\le D(n,j).
\]
If candidate $n$ has a NIP at position $j$, then
\[
 N(n,j)+1\le D(n,j).
\]
\end{theorem}

\begin{proof}
Apply the lowering injection to the NIP occurrences in the prefix. Each image
belongs to an earlier completed candidate and is therefore among the prefix's
DROP occurrences. The source and target occurrence lists have no duplicate
occurrence identifiers, so injectivity gives the first cardinality inequality.

For the stronger result, include the current NIP occurrence $(n,j)$ among the
injection's sources even though it is not yet in the prefix. Its image still
lies in an earlier completed candidate. It is distinct from the images of all
earlier NIPs. There are therefore at least $N(n,j)+1$ earlier DROPs.
\end{proof}

The formal statements have no assumptions about codec injectivity, monotonicity,
surjectivity, a generic enumeration, or a distinctness invariant. Those facts
are proved internally. The prefix theorem remains valid for $j$ beyond the
length of the current word, since invalid positions belong to neither filter.

\subsection{Protected stack tail}

The intended stack invariant is a logical witness of an explicit list of
proved formula heads above an arbitrary opaque initial tail. If its available
height $h$ satisfies
\[
 D(n,j)\le h+N(n,j),
\]
then the strict credit theorem gives $h>0$ immediately before a NIP.
DROP provides one head, NIP consumes one, and the remaining command contracts
are required not to reduce the conservative credit assigned here. AX may be
assigned zero credit in this argument.

That last observation is a proof simplification, not authorization to remove
AX outputs from a machine. Such a machine change could alter completeness or
the numerical calculus and would need its own calibration. Similarly, the
logical height is not a numerical size measure of the encoded stack. Since
$\pair00=0$, zero can encode multiple explicit truth heads above a zero tail.
A witness list, rather than the natural number's magnitude, is the correct
invariant.

\texttt{ConcretePrefixCredit.v} has been compiled with Coq 8.18.0 and recursively
checked with \texttt{coqchk}. The two parser inverses, injection, two credit
theorems, and protected-head corollary all print closed global-context reports.
The delivered reconstruction script records exact source hashes, commands,
compiler output, and kernel-check results. A separate finite parser comparison
binds the conventions to the released source model; it supplements rather
than replaces the unbounded formal proof.

That comparison extracts the actual proved Coq parser to OCaml and matches
every command and field of candidates zero through 999. It also interprets
the literal 4,752-leaf source: 4,000 isolated field-reader cases, 1,000 refills,
1,000 positive continuations, and 1,000 dispatch paths pass, using 19,686,389
source transitions. Complete AX/MP bodies are outside this check. The remaining
unbounded parser/loop bridge is required before the canonical prefix theorem
can be presented as a theorem about every literal machine prefix.

\section{Which 215-state candidate is easier to prove?}

\subsection{Shortest additional proof: CM215}

CM215 is the best current choice for a short \emph{relative} correctness proof.
Its six-row change from K215 admits a whole-tape simulation at identical head
positions, with each old step represented by one or three new steps. The
relation maps old control 176 to new control 52 and otherwise retains controls.
Three simple local tape conditions distinguish the merged roles. Incoming
proposals already contain a compiled bridge to the exact existing dense/Core
machine semantics, so re-proving its numerical calculus, enumerator, compiler,
and blank startup would duplicate established work.\cite{cm215,incoming}

The reusable principle is worth stating with all hypotheses visible.
Here a terminal configuration means one whose partial successor function is
undefined, and terminal status agrees in both directions at every related pair.

\begin{proposition}[Positive-block transport]
For deterministic partial transition systems, suppose an initial relation
preserves terminal observations, preserves a source invariant, and maps each
nonterminal source step to a strictly positive finite target block ending in
the relation. Then the related initial states have equivalent halting
predicates.
\end{proposition}

\begin{proof}
Forward transfer concatenates the simulated blocks. For reverse transfer,
induct on a finite target halting length. The source-terminal case is decided
by terminal agreement. Otherwise the positive target block consumes at least
one target step; determinism makes that block a prefix of any target halting
run, so cancel it and apply induction at the related endpoints. The decrease
would fail for zero-length blocks. This theorem is already checked in the
incoming proposals; it is cited here as an integration tool, not as a new
contribution.
\end{proof}

This does not make CM215's entire proof development independent of K215.
It gives a short new layer over a shared certified base. The final ZFC adapter
must be compiled in the established project before advertising a freshly
rebuilt CM215 end-to-end endpoint. The incoming packets themselves distinguish
that wrapper from their completed local bridge.

\subsection{Simplest scheduler explanation: ZC215}

ZC215 changes the scheduler and startup, retaining the command bodies of its
command-permuted parent. Its clean raw enumeration and lowering injection
make the stack-protection argument particularly transparent. The concrete
credit theorem in this packet now supplies a checked implementation of that
mathematical advantage. But the source-loop and binary-startup connections
are additional integration work; the theorem alone is not the complete
ZC215-to-ZFC endpoint.

The combined GC-ZC215 alternative makes further changes to generalization and
decoder allocation. Those can be interesting for state reduction, but they
introduce additional obligations if proof simplicity alone is the objective.
The four-row FM215 alternative is also attractive for a compact local proof,
and its incoming packet already supplies one. Its remaining repository
semantics adapter should be counted when comparing current formal readiness.

\begin{center}
\begin{tabular}{L{.18\textwidth}L{.35\textwidth}L{.37\textwidth}}
\toprule
Candidate & Reason to choose it & Work still relevant\\
\midrule
Primary K215 & Existing full checked endpoints & Adopt refactors with unchanged
public statements\\
CM215 & Very short exact relative machine proof & Final project endpoint
composition/replay\\
ZC215 & Direct all-prefix credit argument & Literal parser/loop, startup, and
binary endpoint integration\\
FM215 & Small four-row local transformation & Repository representation
adapter and endpoint composition\\
GC-ZC215 & Combined encoding research & Extra generalization and decoder
obligations\\
\bottomrule
\end{tabular}
\end{center}

\section{A complementary, reproducible search for 214 states}

\subsection{Predecessor-specific elimination}

Previous searches explored both fixed-pair merges and cutpoint rewrites. In
particular, the newest incoming proposal contains a two-row cutpoint search
for helper 176, so deleting a state without one uniform state-map image is
not a new idea. The present contribution systematically applies per-incoming-row
choices to every eligible deletion target in 21 released tables, together with
a separate state-40 provider experiment. Its composability restrictions differ
from the earlier cutpoint search; neither grammar contains the other in general.

Choose a noninitial control $d$ to delete. For every external transition row
that enters $d$, independently choose a replacement write, direction, and
retained destination or halt. Different
predecessors may route to different continuations. Thus successful elimination
need not identify $d$ with one common surviving state.

For each proposed replacement row and each admitted local tape window, the
checker follows the old and edited computations until they meet at exactly
the same surviving control, head displacement, and full local tape contents.
Neither path may leave the tracked window before the certified meeting. Thus
all cells outside that window are framed unchanged. Nontrivial blocks have
positive duration. Other simultaneously changed rows are excluded from the
macro interior; repeated uses of the currently checked row are handled with
the proposed replacement. These restrictions make independently accepted
incoming-row implementations composable.

The input windows are computed by a finite abstract reachability closure.
The implemented closure starts from every local window at the declared start
control; when the head moves, an unobserved boundary bit is conservatively
allowed either value. Thus it overapproximates starts on arbitrary tapes,
including blank input. This stronger invariant supports sound positive
certificates, but can reject a valid blank-only optimization because an
admitted window is actually unreachable from the blank configuration.

\subsection{Completed scope and outcome}

\SearchMeasurements

The principal run uses radius five and a 96-step local trace cap on 21 released
215-state tables. A second run tests the primary at radius seven and a larger
192-step cap. There is also a targeted search that changes a provider row while
trying to bypass state 40. Positive controls and exact input hashes accompany
the code and receipts.

The 21 inputs have distinct rooted transition graphs, and remain 21 distinct
classes when global reflection is allowed. All 215 controls are graph-reachable
in each table. The dominant rejection category is leaving the finite window
(6,337,648 local trials). These counts describe failed sufficient certificates,
not semantic counterexamples to whole machines. The local-join transformation
theorem in the search README is an ordinary mathematical argument; it is not
advertised as a new kernel-checked theorem.

The quantities called action trials are local row proposals, not numbers of
distinct complete Turing machines. The broad run does not enumerate all
combinations of successful incoming-row choices when one required predecessor
has no accepted implementation. This early rejection is sufficient for the
specific independently composable grammar but says nothing about mutually
dependent coordinated rewrites outside it.

\subsection{What the closest obstruction teaches}

On the primary table, state 40 provides a useful near miss: two of
its three external incoming rows can be rerouted, while the third remains
unresolved. States 40 and 54 jointly scan left for the first double zero. The
two roles remember whether a single zero has just been seen. That explains
why treating them as interchangeable decoder continuations loses necessary
information.

This observation is not a general lower bound for a one-tape Turing machine.
A different encoding, a temporary tape mark, a stronger reachable-state
invariant, or a simulation relation with shifted comparison points might
store that information elsewhere. The failed provider search records one
finite attempt to do so with existing rows; it does not rule out those broader
possibilities.

An independent dictionary-tape replayer reduces the failure at predecessor
\texttt{54/0} to three admitted radius-five windows:
\texttt{00001010000}, \texttt{00101010000}, and \texttt{00000010100}, with
the head on the sixth bit. Of 860 possible replacement actions, three pass
the first window, two pass the first two, and none passes all three under the
specified bounds. These are abstract-invariant contexts; their occurrence on
the blank run is not claimed. A separate positive control deletes a live shared
state using two different providers and is independently replayed on all 64
nonvacuous contexts.

The search treats fuel exhaustion, leaving the tracked window, and failure to
find a bounded exact join as rejection by the template. They are not all
counterexamples to the proposed whole machine. No claim of minimality follows
from a large action-trial count.

\section{Integration plan and further research}

\subsection{Recommended incorporation sequence}

\begin{enumerate}
\item\textbf{Adopt the normalized Rocq source layer.} Preserve the literal
source programs and existing theorem interface. Add the reusable checker and
replace the large offset-stage source proof by its 1,426-entry one-sided
instance. Replay under the maintained Rocq toolchain and audit the full endpoint.
\item\textbf{Choose how to maintain Lean correspondence data.} The incoming
normalized proof is already a good intermediate state. The automatic version
generates its materialized certificate from the source trees. Keep that staged
architecture and measure its cost against the existing materialized version
before choosing the production default.
\item\textbf{Add CM215 through the existing exact transfer.} Reuse the common
K215 calibration and complete the final checked project composition. Avoid
duplicating source or bootstrap proofs in each small table variant.
\item\textbf{Integrate the concrete ZC215 codec and credit theorem.} Connect
the implementation's command/refill semantics to the exact parser, then use
the protected-head corollary to discharge stack safety. Keep the arbitrary
tail as an explicit logical witness.
\item\textbf{Continue state synthesis with proof-producing interfaces.} A
future successful 214 table should come with a literal state count, exact
initial configuration, a total runtime relation, and positive progress or a
well-founded stuttering measure. Bounded agreement alone is insufficient.
\end{enumerate}

\subsection{A concrete next attempt at 214 states}

The most actionable new synthesis direction is to weaken the local physical
register-file contract. The existing \texttt{Runtime230Core.v} increment box
fixes its output to an exact file of the form
\texttt{sepsR (incr r vs ++ [1])}. Its precondition supplies at least ten file
entries while the simulated program accesses only the first nine logical
registers. At the higher layer, \texttt{Sim230.v} already hides the spare
extension existentially. Thus the proposed change is to the stronger local
box specification, not the invention of a new global abstraction.\cite{runtime}

Preserve the first nine registers and the logical continuation, but permit
any nonempty valid spare suffix and an allowed gap on return:
\[
 \operatorname{enc}(p,\mathbf r,g,s)
 \ \longrightarrow^{+}\ 
 \operatorname{enc}(p',\mathbf r',g',s'),\qquad s'\ne[].
\]
Then resynthesize the register excursion, spare append, and return-to-dispatch
macro jointly. This permits changes to old spare values and the number of
spare blocks, beyond choosing whether a newly appended block contains zero or
one. Exact whole-tape joins exclude these changes before their abstract
semantic validity can be considered.

The new macro would still need a universal proof preserving the nine logical
registers, well-formed nonempty suffix, permitted gap, every shared decoder
use, positive progress, and exact blank startup. No machine realizing a saving
under this relaxed contract has yet been constructed. The search packet gives
the precise existing formal-source locations behind this proposal.

There is also a useful dependency cleanup. The file \texttt{Prog215.v}
imports \texttt{K230.Prog230} for its encoded
edge constructors and \texttt{cv}, pulling in an unrelated 5,124-node source
table and well-formedness certificate. Extract those shared constructors into
a small generic module imported by both programs. This simplification is
proposed; its runtime effect has not been measured.

\subsection{Specific further questions}

\begin{enumerate}
\item Can automatic correspondence inference be specialized to verified
annotated source DAGs, so repeated subtrees share checkpoint counts and
segment summaries during kernel reduction? This would address computation
cost while retaining the no-external-correspondence interface.
\item Can the normalization-aware relation be generalized from trailing-zero
aliases to equality of the complete consumed primitive observation? That
larger quotient can merge pointers with unused suffixes, but it requires
careful successor and halt preservation.
\item Can the current direct halting-witness transport in the Rocq checker
be factored into a generic segment relation theorem matching Lean's
\texttt{Bisim} interface? The present implementation prioritizes a small,
checked port over a broad library refactor.
\item Can the literal ZC215 register parser be related to the concrete codec
by a short symbolic block proof, so the new credit theorem removes the older
debt-accounting argument from an entire candidate endpoint?
\item Does a table with a slightly larger source but a simpler bootstrap admit
a smaller total binary allocation? Source-DAG size, occurrence count, and
working-state count should be optimized jointly, not used interchangeably.
\item Can the double-zero scanner's memory be moved into a temporary local
marker that is restored at a different comparison boundary, eliminating a
working control without adding a provider? The present exact local joins
do not settle that representation change.
\item Can a reachable-phase certificate permit additional junk formula or
stack configurations while preserving both directions of the consistency
calibration? Equiconsistency needs soundness and completeness, even when the
intermediate syntax has no intended set-theoretic interpretation.
\item After the source and local machine transport layers are compact, what
is the smallest additional formal development needed to obtain an internally
encoded PA or ZFC proof for the exact 215-state table? This remains distinct
from the external proof-assistant theorem.
\end{enumerate}

\section{Artifact and verification guide}

The archive is designed as an additive research packet. It does not modify or
replace the repository checkout used for comparison. Its main components are
\texttt{lean\_auto/}, \texttt{rocq\_normalized/}, \texttt{rocq\_auto/}, \texttt{zc\_credit/},
\texttt{search214/}, \texttt{evidence/}, this article, and its editable source.
Each component's README gives a portable rebuild command and states the
boundary between its new theorem and inherited project dependencies.

The top-level manifest binds the delivered files. Hash checks establish artifact
identity; they do not replace kernel checking. The packet retains compiler
and kernel logs for the newly checked theorems and contains the source needed
to reproduce those checks. It does not distribute toolchain binaries or
compiled proof objects. Search receipts record their input tables, abstraction
bounds, and unsuccessful outcomes.

The report's new formal claims are restricted to the recorded source and
scheduler theorems. The original machine-to-ZFC theorem, the prior CM215
bridges, the shared-tree representation, and the normalized Lean checker are
explicit inherited dependencies. No absolute consistency assertion, 214-state
success, or universal state lower bound is claimed.

\begingroup\small
\setlength{\parskip}{2pt}
\begin{thebibliography}{9}
\setlength{\itemsep}{2pt}\setlength{\parsep}{0pt}
\bibitem{verification} ProveIt, \emph{K215 formal verification record},
commit \texttt{53269bcab1ef}.
\href{https://github.com/VladimirReshetnikov/ProveIt/blob/53269bcab1efc0f82d3022645bd7cf0c22d84959/Computability/BusyBeaver/Research/ZFC/k225-formal/verification/k215/README.md}{Pinned verification README}.
\bibitem{phase4} ProveIt, \emph{Shared implication lemmas and symbolic execution},
phase-four simplification record.
\href{https://github.com/VladimirReshetnikov/ProveIt/blob/53269bcab1efc0f82d3022645bd7cf0c22d84959/Computability/BusyBeaver/Research/ZFC/k225-formal/verification/k215/simplification/phase4/README.md}{Pinned phase-four README}.
\bibitem{incoming} ProveIt, seven pending research archives in
\texttt{docs/incoming}, inspected at the pinned commit.
\href{https://github.com/VladimirReshetnikov/ProveIt/tree/53269bcab1efc0f82d3022645bd7cf0c22d84959/docs/incoming}{Pinned incoming directory}.
Exact archive names and SHA-256 digests accompany this packet.
\bibitem{cm215} ProveIt, \emph{CM215: conditional origin markers and ZFC consistency}.
\href{https://github.com/VladimirReshetnikov/ProveIt/blob/53269bcab1efc0f82d3022645bd7cf0c22d84959/Computability/BusyBeaver/Research/ZFC/candidate-k215/alternatives/conditional-marker/README.md}{Pinned CM215 construction and ordinary proof}.
Formal follow-up status is in the incoming archives.
\bibitem{zc215} ProveIt, \emph{ZC215: zero-credit enumeration}, including
\texttt{proofs/tag\_lowering\_credit.md}.
\href{https://github.com/VladimirReshetnikov/ProveIt/blob/53269bcab1efc0f82d3022645bd7cf0c22d84959/Computability/BusyBeaver/Research/ZFC/candidate-k215/alternatives/zero-credit/README.md}{Pinned ZC215 README}.
\bibitem{pathsemantics} ProveIt, \emph{RegisterPath, RMEquiv, and RMBisim},
maintained Lean source semantics.
\href{https://github.com/VladimirReshetnikov/ProveIt/tree/53269bcab1efc0f82d3022645bd7cf0c22d84959/Computability/BusyBeaver/Lean/BusyBeaver/ZFC}{Pinned Lean module directory}.
The files preserve attribution to Wade's compiler framework and the
Mira-Cult/Jeff verification on which their semantics are based.
\bibitem{runtime} ProveIt, \emph{Runtime230Core and Sim230}, selected runtime
contracts for K215.
\href{https://github.com/VladimirReshetnikov/ProveIt/blob/53269bcab1efc0f82d3022645bd7cf0c22d84959/Computability/BusyBeaver/Research/ZFC/k225-formal/rocq/rocq-extension/coq/215/Runtime230Core.v}{Pinned local runtime contracts}.
The existential representation is in the adjacent \texttt{Sim230.v}.
\end{thebibliography}
\endgroup
\end{document}
